
\begin{theorem}[Theorem 20.1]\label{mcrt-auto-theorem-20-1}\dcref{20.1}\source{matsumura-commutative-ring-theory:page-0176}\uses{mcrt-auto-theorem-13-9}
\begin{mathematicalrecord}
A Noetherian integral domain A is a UFD if and only if
every height 1 prime ideal is principal.
\end{mathematicalrecord}
\end{theorem}
\begin{proof}
\begin{mathematicalrecord}
Suppose that A is a UFD and that P is a height 1
prime ideal. Take any non-zero aeP, and express a as a product of prime
elements, a = nq. Then at least one of the ni belongs to P; if 7liEP then
(xi) c P, but (rq) is a non-zero prime ideal and ht P = 1, hence P = (xi).
Proof of ?iLf?. Since A is Noetherian, every element aEA which is neither
0 nor a unit can be written as a product of finitely many irreducibles.
Hence it will be enough to prove that an irreducible element a is a prime
element. Let P be a minimal prime divisor of (a); then by the principal
ideal theorem (Theorem 13.9, ht P= 1, so that by assumption we can
write P =(b). Thus a = bc, and since a is irreducible, c is a unit, so that
(a) = (b) = P, and a is a prime element.    H
\end{mathematicalrecord}
\end{proof}


\begin{theorem}[Theorem 20.2]\label{mcrt-auto-theorem-20-2}\dcref{20.2}\source{matsumura-commutative-ring-theory:page-0177}\uses{}
\begin{mathematicalrecord}
Let A be a Noetherian integral domain, r a set of prime
elements of A, and let S be the multiplicative set generated by f. If A, is
a UFD then so is A.
\end{mathematicalrecord}
\end{theorem}
\begin{proof}
\begin{mathematicalrecord}
Let P be a height 1 prime ideal of A. If PnS # 0 then P contains
an element nil-, and since nA is a non-zero prime ideal we have P = nA.
If PnS = 0 then PA, is a height 1 prime ideal of A,, so that PA, = aA,
for some aEP. Among all such a choose one such that aA is maximal;
then a is not divisible by any nil-. Now if XEP we have S.X=ay for
some SES and YEA. Let s = or... x I with n,gT; then a$z,A, so that
y~qA, and an induction on r shows that y~sA, so that xEaA. Hence
P=aA.     n

Lemma 1. Let A be an integral domain, and a an ideal of A such that
 a@A?2:Anf1;          then a is principal.
Proof. Fix the basis e,, . , e, of A??,          and viewing a @ A? c A @A?, fix
fO,. . . , f, such that ,fO is a basis of A and fi,. . , f,, a basis of A?. Then
the isomorphism          q:A?+ ? -a@       A? can be given in the form q(ei) =
C;=o~iJj        Write di for the (i,O)th cofactor of the matrix (aij), and d
for the determinant, so that, since cp is injective, rl # 0, and 1 Uiodi = d,
C~ijdi =O if j # 0. H ence if we set eb = C;d,e, we have cp(eb) = dfo.
Moreover,       since the image of cp includes Jr,.. .,fn, there exist e;,...?
e;EA ?+? such that cp(e>) =fj. Now define a matrix (c,) by e; = x;=ocjkek
for j = 0,. . . , n (so c,,~ = C&J.Then we have
                         /d    0   . ..      O\

        (Cjk)(aij)   =    (f   ?   ..        O

                         i 0   0        .i   1?
SO that by comparing the determinants of both sides we get det tcjk)       =   I?
Therefore eb,. . , e; is another basis of A?+l, and afo = cp(Aeb) = dAf,,
  so that a = dA. n
     Let K be the field of fractions of the integral domain A; for a finite
  A-module M, the dimension of M aA K as a vector space over K is called
  the rank of M. A torsion-free finite A-module of rank 1 is isomorphic to
  an ideal of A. Lemma 1 can be formulated as saying that for an integral
  domain A, a stably free rank 1 module is free (see Ex. 19.2). The elementary
  proof given above is taken from a lecture by M. Narita in 1971.
\end{mathematicalrecord}
\end{proof}


\begin{theorem}[Theorem 20.3]\label{mcrt-auto-theorem-20-3}\dcref{20.3}\source{matsumura-commutative-ring-theory:page-0178}\uses{mcrt-auto-theorem-7-12}
\begin{mathematicalrecord}
Auslander and Buchsbaum [3]). A regular local ring is a
  UFD.
\end{mathematicalrecord}
\end{theorem}
\begin{proof}
\begin{mathematicalrecord}
Let (A,m) be a regular local ring: the proof works by induction
  on dim A. If dim A = 0 then A is a field and therefore (trivially) a UFD.
  If dim A = 1 then A is a DVR, and therefore a UFD. We suppose that
  dim A > 1 and choose xEm - m2; then since xA is a prime ideal, applying
  Theorem 2 to I- = {x}, we need only show that A, is a UFD (where
  A, = A[x- ?1 is as on p. 22). Let P be a height 1 prime ideal of A, and
  set p = Pn A; we have P = pA,. Since A is a regular local ring, the
  A-module p has an FFR, so that the AX-module P has an FFR. For any
  prime ideal Q of A,, the ring (A,), = A,,, is a regular local ring of
  dimension less than that of A, so by induction is a UFD. Thus P, is free
  as an (A&-module,     so that by Theorem 7.12, the AX-module P is projective;
, hence by Ex. 19.2, P is stably free, and therefore by the previous lemma,
  P is a principal ideal of A,.    n
     The above proof is due to Kaplansky. Instead of our Lemma 1, he used
  the following more general proposition, which he had previously proved:
  if A is an integral domain, and Zi,Ji are ideals of A for 1 d i d r such that
  a=, Ii N @= 1 Ji, then I,. . I, N J,. . . J,. This is an interesting property
  of ideals, and we have given a proof in Appendix C.
\end{mathematicalrecord}
\end{proof}


\begin{theorem}[Theorem 20.4]\label{mcrt-auto-theorem-20-4}\dcref{20.4}\source{matsumura-commutative-ring-theory:page-0178}\uses{mcrt-auto-theorem-7-12}
\begin{mathematicalrecord}
Let A be a Noetherian integral domain. Then if any finite
  A-module has an FFR, A is a UFD.
\end{mathematicalrecord}
\end{theorem}
\begin{proof}
\begin{mathematicalrecord}
By Ex. 19.4, A is a regular ring. Let P be a height 1 prime ideal
  of A. Then A,,, is a regular local ring for any mESpec A, so by the previous
  theorem, the ideal P,,, is principal, and is therefore a free A,,,-module. Hence
  by Theorem 7.12, P is projective. Therefore by Ex. 19.2, P is stably free,
  and so by Lemma 1 is principal.         n
     Let A be an integral domain; for any two non-zero elements a, bEA,
  the notion of greatest common divisor (g.c.d.) and least common multiple
  (1.c.m.) are defined as in the ring of integers. That is, d is a g.c.d. of a and
  b if d divides both a and b, and any element x dividing both a and b
164      Regular rings

divides d; and e is an 1.c.m. of a and b if e is divisible by both a and b,
and any y divisible by a and b is divisible by e; this condition is equivalent
to (e) = (a) n(b).
Lemma 2. If an 1.c.m. of a and b exists then so does a g.c.d.
Proof. If (a)n(b) = (e) then there exists d such that ab = ed. From ee(a)
we get bQd) and similarly a@d), so that (a, b) c (d). Now if x is a common
divisor of a and b then a = xt and b = xs, so that xst is a common multiple
of a, b, and is hence divisible by e. Then from ed = ab = x.xst we get that
d is divisible by x. Therefore, d is a g.c.d. of a and b. n

Remark 1. If A is a Noetherian    integral domain which is not a UFD then
A has an irreducible element a which is not prime. If xyE(a) but x$(a),
 y+(a) then the only common divisors of a and x are units, so that 1 is a
g.c.d. of a and x. However, xye(a)n(x),          but xy$(ax), so that (a)n
(x) # (ax), and there does not exist any 1.c.m. of a and x. Thus the converse
of Lemma2 does not hold in general.
Remark 2. If A is a UFD then an intersection        of an arbitrary collection
of principal ideals is again principal (we include (0)). Indeed, if n ill aiA # 0,
then factorise each a, as a product of primes:



with ui units, and p, prime elements such that pJ # p,A for x # 8.
Then r) a,A = dA, where d = ~~~~~~~~~~~~~~~~~~
                                             (We could even allow the a, to
be elements of the field of fractions of A.)
\end{mathematicalrecord}
\end{proof}


\begin{theorem}[Theorem 20.5]\label{mcrt-auto-theorem-20-5}\dcref{20.5}\source{matsumura-commutative-ring-theory:page-0179}\uses{}
\begin{mathematicalrecord}
An integral domain A is a UFD if and only if the ascending
chain condition holds for principal ideals, and any two elements of A have
an 1.c.m.
ProoJ: The ?only if? is already known, and we prove the ?if?. From the first
condition it follows that every element which is neither 0 nor a unit can
be written as a product of a finite number of irreducible elements, so that
we need only prove that an irreducible element is prime. Let a be an
irreducible element, and let xyE(a) and x$(a). By assumption we can
write (a)n(x) = (2); now 1 is a g.c.d. of a and x, so that one sees from the
\end{mathematicalrecord}
\end{theorem}
\begin{proof}
\begin{mathematicalrecord}
that YE(a). Therefore (a) is prime.    n
\end{mathematicalrecord}
\end{proof}


\begin{theorem}[Theorem 20.6]\label{mcrt-auto-theorem-20-6}\dcref{20.6}\source{matsumura-commutative-ring-theory:page-0179}\uses{mcrt-auto-theorem-11-3,mcrt-auto-theorem-12-3,mcrt-auto-theorem-7-11}
\begin{mathematicalrecord}
Let A be a regular ring and u, UEA. Then uAnuA is a
projective ideal.
\end{mathematicalrecord}
\end{theorem}
\begin{proof}
\begin{mathematicalrecord}
A,,, is a UFD for every maximal ideal m, so that (uAnuA)A,, =
uA,nvA,     is a principal ideal, and hence a free module. n
?Remark.     There are examples where A is a UFD but A[X] is not.
     i It is easy to see that a UFD is a Krull ring. For any Krull ring A we
     : can define the divisor class group of A, which should be thought of as a
     ?Ineasure of the extent to which A fails to be a UFD. We can give the
     ,fdeIinition in simple terms as follows: let 9 be the set of height 1 prime
       ideals of the Krull ring A, and D(A) the free Abelian group on 9. That
     ?is, D(A) consists of formal sums 1 ptBnp?p (with n,EZ and all but finitely
     ?many n, = 0), with addition defined by
                (Cn;p)+(Cnb.P)=C(n,+n6)p.
          K be the field of fractions of A, and K* the multiplicative   group of
               o elements of K, and for EK* set div(a) = CPE,~uP(u).p, where up
                ormalised additive valuation of K corresponding to p. Then
               = div(a) + div (b), so that div is a homomorphism from K* to D(A).
          write F(A) for the image of K*; this is a subgroup of D(A), so that
          can define C(A) = D(A)/F(A) to be the divisor class group of A.
             US~Y, if A is a UFD then each PEP is principal, and if p = aA then
             element of D(A) we have p = div(a), so that C(A) = 0. Conversely,
            ) = 0 then each PEP is a principal ideal, and putting this together
           he corollary of Theorem 12.3, one sees easily that A is a UFD. Hence
                A is a UFD~C(A)       = 0.
Now let A be any ring, and M a finite projective A-module. For each
PtsSpecA, the localisation M, is a free module over A,, and we write
n(P) for its rank. Then n is a function on Spec A, and is constant on every
connected component (since n(P) = n(Q) if P 2 Q). This function n is called
the rank of M. If the rank is a constant r over the whole of SpecA then
we say that M is a projective module of rank r. We write Pit(A) for the
set of isomorphism    classes of finite projective A-modules of rank 1; cl(M)
denotes the isomorphism class of M. If M and N are finite projective rank
1 module then so is MOAN; this is clear on taking localisations. Thus
we can define a sum in Pit(A) by setting
       cl(M)+cl(N)=cl(M@N).
We set M? = Hom,(M,      A), and define cp:M 0 M* -A         by
         44x   mi   0   fi)   =   C   h(F);

then cp is an isomorphism (taking localisations and using the corollary to
Theorem 7.11 reduces to the case M = A, which is clear). Hence cl(M*) =
 -cl(M),    and Pit(A) becomes an Abelian group, called the Picard group
of A. If A is local then Pit(A) = 0.
    If A is an integral domain with field of fractions K, then MC,, = M 0 K,
so that the rank we have just defined coincides with the earlier definition
(after Lemma 1). If M is a finite projective rank 1 module, then since M
is torsion-free    we have M c MC,, N K, so that M is isomorphic as an
A-module to a fractional ideal; for fractional ideals, by Theorem 11.3,
projective and invertible are equivalent conditions, so that for an integral
domain A, we can consider Pit(A) as a quotient of the group of invertible
fractional ideals under multiplication.   A fractional ideal I is isomorphic
to A as an A-module precisely when I is principal, so that




   Suppose in addition that A is a Krull ring. Then we can view Pit(A)
as a subgroup of C(A). To prove this, for p~;/p and I a fractional ideal.
set
         v,(Z) = min {v,(x)Ix~l};
this is zero for all but finitely many PEP (check this!), so that we can set
         div (I) = 1 v,,(l).p~D(A).
                   PEP
For a principal ideal I = zA we have div(I) = div(cx). One sees easily that
div(Z1?) = div(1) + div(I?), and that div(A) = 0, so that if I is invertible,
div (I) = - div (I - ?).
For invertible I we have (I-?))?          = I: indeed, I c(Z-?)-I       from the
  definition, and I = Z.A 2 Z(Z-?(I-?)-?)          1 (I-?)-?. If Z is invertible and
  ~v(Z)=Othendiv(Z-?)=O,sothatZcA,Z-?cA;henceAc(Z-?)-1=Z,
  and Z = A. It follows that if I, I? are invertible, div(Z) = div(Z?) implies Z = I?.
  Thus we can view the group of invertible fractional ideals as a subgroup of
  D(A), and Pit(A) as a subgroup of C(A).
     If A is a regular ring then as we have seen, p~9? is a locally free module,
  and so is invertible. Clearly from the definition, div(p) = p. Hence, in the
~ case of a regular ring, D(A) is identified with the group of invertible
1,fractional ideals, and C(A) coincides,with Pit(A).
     The notions of D(A) and Pit(A) originally arise in algebraic geometry.
,: Let V be an algebraic variety, supposed to be irreducible and normal. We
1,write 9 for the set of irreducible codimension 1 subvarieties of V, and
Sdefine the group of divisors D(V) of I/ to be the free Abelian group on 9?;
; a divisor (or Weil divisor) is an element of D(V). Corresponding to a
       onal function f on I/ and an element WEP, let z+(f) denote the
       er of zero off along W, or minus the order of the pole if f has a pole
       ng W. Write div (f) = ~,,9uw(f). Wfor the divisor of f on I? (or just
         For WEP, the local ring CO, of W on I/ is a DVR of the function
         of V, and ow is the corresponding valuation. We say that two divisors
            D(V) are linearly equivalent if their difference M - N is the divisor
              ction, and write M - N. The quotient group of D(V) by - , that
          quotient by the subgroup of divisors of functions, is the divisor class
            of I/ (up to linear equivalence), and we write C(V) for this. (In
         ion to linear equivalence one also considers other equivalence
      ations with certain geometric significance (algebraic equivalence,
      merical equivalence,. . . ), and d ivisor class groups, quotients of D(V) by
           rresponding subgroups.)
           ivisor M on I/ is said to be a Cartier divisor if it is the divisor of
         ction in a neighbourhood of every point of I/. From a Cartier divisor
          onstructs a line bundle over V, and two Cartier divisors give rise to
      morphic line bundles if and only if they are linearly equivalent. Cartier
         ors form a subgroup of D(V), and their class group up to linear
          alence is written Pit(V); this can also be considered as the group of
      morphism classes of line bundles over V (with group law defined by
       sor product). If T/ is smooth then (by Theorem 3) there is no distinction
     tween Cartier and Weil divisors, and C(V) = Pit (V).
         e reader familiar with algebraic geometry will know that the divisor
      ss group and Picard group of a Krull ring are an exact translation of
        corresponding notions in algebraic geometry. If I/ is an affine variety,
     th coordinate ring k[ I?] = A then C( I?) = C(A) and Pit (I?) = Pit (A). In
this case, to say that A is a UFD expresses the fact that every codimension
  1 subvariety of I/ can be defined as the intersection of I/ with a hypersurface.
 If I/ c P? is a projective algebraic variety, defined by a prime ideal
 Zck[X,,...,      X,], and we set A = k[X]/I = k[to,. . . , t,,] (with ti the
 class of Xi) then A is the so-called homogeneous coordinate ring of I/. If
 A is integrally closed we say that V is projectively normal (also
 arithmetically normal). This condition is stronger than saying that V is
 normal (the local ring of any point of I/ is normal). If A is a UFD then
 every codimension 1 subvariety of Vcan be given as the intersection in
 P? of V with a hypersurface. Let m = (Co,. . . ,&,) be the homogeneous
 maximal ideal of A, and write R = A,,, for the localisation. The above
 statement holds if we just assume that R is a UFD; see Ex. 20.6. All the
 information about V is contained in the local ring R.
     Thus C(A), Pit(A) and the UFD condition are notions with important
 geometrical meaning, and methods of algebraic geometry can also be used
 in their study. For example, in this way Grothendieck [G5] was able to
 prove the following theorem conjectured by Samuel: let R be a regular
 local ring, P a prime ideal generated by an R-sequence, and set A = R/P;
 if A, is a UFD for every p~Spec A with ht p < 3 then A is a UFD.
     We do not have the space to discuss C(A) and Pit(A) in detail, and we
just mention the following two theorems as examples:
     (1) If A is a Krull ring then C(A) 2: C(A[X]).
     This generalises the well-known theorem (see Ex. 20.2) that if A is a UFD
then so is A[X].
     (2) If A is a regular ring then C(A) N C(A[XJ).
     This generalises Theorem 8.
     Finally we give an example. Let k be a field of characteristic 0, and set
 A = k[X, Y, Z]/(Zn - X Y) = k[x, y, z] for some n > 1. Then A/(z, x) N
 k[X, Y, 21/(X, Z) N k[ Y], so that p = (x, z) is a height 1 prime ideal of A.
 In D(A) we have np = div(x), and it can be proved that C(A) z Z/d (see
 [S2], p. 58). The relation xy = z? shows that A is not a UFD.
     For those wishing to know more about UFDs, consult [K], [S21 and
IFI.
\end{mathematicalrecord}
\end{proof}


\begin{theorem}[Theorem 21.1]\label{mcrt-auto-theorem-21-1}\dcref{21.1}\source{matsumura-commutative-ring-theory:page-0185}\uses{}
\begin{mathematicalrecord}
Let (A,m, k) be a Noetherian            local ring, and A its
completion.
  (i) E&A) = ~~(2) for all p 3 0.
  (ii) el(A) 3 em bdimA-dimA.
  (iii) If R is a regular local ring, a an ideal of R and A N R/a, then
          p(a) = dim R - emb dim A + cl(A).
\end{mathematicalrecord}
\end{theorem}
\begin{proof}
\begin{mathematicalrecord}
(i) is clear from the fact that a minimal basis of m is a minimal
basis of mA^, so that applying OaA^ to the complex E, made from A
gives that made from A. Then since A is A-flat, H,(E.) @ A^ = H,(E. 0 A^),
and mH,(E.) = 0 gives H,(E.) @ A^= H,(E.).
   (ii) If A is a quotient of a regular local ring, then as we have seen above,
there exists a regular ring (R, n) such that A = R/a with a c n2, so that
el(A) = p(a) > ht a = dim R - dim A = emb dim A - dim A,           where    the
equality for ht a comes from Theorem 17.4, (i). Now A itself is not necessarily
a quotient of a regular local ring, but we will prove later (see 429) that A
always is, and we admit this in the section. Having said this, the two sides
of (ii) are unaltered on replacing A by A, and the inequality holds for A.
   (iii) Set n = rad(R). If a c n2 then, as we have seen above, p(a) = &i(A),
and dim R = embdim A, so that we are done. If a$ n2, take xEa - nZ;
we pass to R/xR and a/xR, each of dim R and p(a) decreases by 1, SO
    nduction completes the proof. n
    ition. A Noetherian local ring A is a complete intersection ring
    eviated to c.i. ring) if cl(A) = emb dim A - dim A.

heorem 21.2. Let A be a Noetherian        local ring.

 (ii) Let A be a c.i. ring and R a regular local ring such that A = R/a;
hen a is generated by an R-sequence. Conversely, if a is an ideal generated
   an R-sequence then R/a is a c.i. ring.
  iii) A necessary and sufficient condition for A to be a c.i. ring is that
     completion A^ should be a quotient of a complete regular local ring
        n ideal generated by an R-sequence.

      By Theorem 1, (iii), p(a) = dim R -embdim          A+ cl(A), and by
    rem 17.4, (i), ht a = dim R - dim A, so that A is a c.i. ring is equivalent
                         Theorem 17.4, (iii), this is equivalent to a being

        e sufficiency is clear from (i) and (ii). Necessity follows from the
    hat A^is a quotient of a complete regular local ring (see ?29), together


  eorem21.3. A c.i. ring is Gorenstein.
  oaf. If A is c.i. then so is A, and if A^is Gorenstein then so is A, so that we
    assume that A is complete. Then we can write A = R/a, where R is
   gular local ring and a is an ideal generated by a regular sequence. Since
  1s Gorenstein, A is also by Ex. 18.1. w
                                  g chain of implications for Noetherian local

       regular * c.i. + Gorenstein s CM.
      A be a c.i. ring, and p a prime ideal of A. If A is of the form
                xI), where R is regular and .x1,. . . , x, is an R-sequence, then
         can be written A, = R,/(x,, . . . , x,), where R, is regular and
         is an R,-sequence, it follows that A, is again a c.i. ring. The
          f deciding whether A, is still a c.i, ring even if A is not a quotient
        ar local ring remained unsolved for some time, but was answered
         ely by Avramov [l], making use of AndrC?s homology theory
        1. This theory defines homology and cohomology                     groups
           ) and H?(A, B, M)for n 3 0 associated with a ring A, an A-algebra
           module M. The definition is complicated, but in any case these
    B-modules having various nice functorial properties. If A is a
Noetherian     local ring with residue lield k then
           A is regular-H&          k, k) = 0,
and
           A is c.i.-H,(A,     k, k) = 0;
for n 2 3 the statements H,(A, k, k) = 0 and H,(A, k, k) = 0 are equivalent.
Thus Andrk homology is particularly relevant to the study of regular and
c.i. rings.
\end{mathematicalrecord}
\end{proof}
